\begin{textsample}{POS Dim 1 – vem\_ed\_03 – Score 8.00 – vem\_ed\_03/t009.txt}  \label{ex:f1_pos_158}
PROCURA-SE Quer desenvolver um Projeto Colaborativo Internacional (PCI)?
Parceiros em Instituições de Ensino estrangeiras buscam pares nas Fatecs em diversas áreas, como: \textbf{Competências} interculturais Gerentes internacionais Mineração e geoprocessamento Mapa de riscos Saúde Meio Ambiente e Sustentabilidade Artes Álgebra e Cálculo História e Ciência Política Caso tenha interesse, basta escrever para cesu.pci@cps.sp.gov.br e juntar-se às nossas equipes no Teams.
QUEBRA-GELO Os projetos \textbf{colaborativos} internacionais desenvolvem \textbf{competências} interpessoais, linguísticas, culturais e digitais, imprescindíveis para os tecnólogos no mundo globalizado.
A Revista do Centro Paula Souza fez uma reportagem sobre os PCIs das Fatecs na edição 76, \textbf{publicada} no fim de agosto.
Esse reconhecimento traz muita satisfação à nossa equipe.
Seguindo nos relatos de sucesso sobre \textbf{intercâmbios} virtuais, esta edição de VEm com PCI apresenta, na seção"Quem \textbf{é} quem", Stephanie Doscher, diretora do escritório de iniciativas de aprendizagem global da Florida International University (FIU).
Desde 2018, a instituição já desenvolveu 18 PCIs com seis Fatecs e, no primeiro semestre deste ano,"Education and Leadership", \textbf{que} envolveu pós-graduandos da FIU e gestores de 15 Fatecs.
A internacionalização \textbf{é} um dos caminhos para a formação dos \textbf{futuros} tecnólogos.
Engajada nessa \textbf{proposta}, a Fatec Indaiatuba fez uma recepção virtual aos estudantes, em 29 de agosto, com palestrantes de universidades da Holanda, dos EUA e da Argentina, com as quais a unidade desenvolve PCIs.
Um exemplo inovador de boas \textbf{práticas}.
Criar, desenvolver e \textbf{expandir} redes de \textbf{intercâmbios} virtuais: eis um dos principais desafios da internacionalização do ensino \textbf{superior}.
Esse \textbf{foi} um dos temas \textbf{que} debati no principal evento sobre \textbf{intercâmbios} virtuais do mundo: International Virtual Exchange Conference (IVEC 2020), realizado de 14 a 16 de setembro.
A tropicalização das experiências de internacionalização em \textbf{casa} \textbf{foi} o assunto \textbf{que} apresentei no Congresso Internacional de Educação e Tecnologias: Encontro de Pesquisadores em Educação a Distância (CIET: EnPED), \textbf{que} ocorreu na última semana de agosto.
Boa leitura!

% matched lemmas: casa, colaborativo, competência, expandir, futuro, intercâmbio, proposta, prático, publicar, que, ser, superior
\end{textsample}
